\documentclass[10pt,letterpaper]{article}
\usepackage[T1]{fontenc}
\usepackage[utf8]{inputenc}
\usepackage[spanish,es-tabla,es-noshorthands]{babel}
\usepackage{geometry}
\geometry{margin=0.68in, top=0.85in, bottom=0.78in, headheight=36pt, headsep=14pt}
\usepackage{amsmath,amssymb}
\usepackage{xcolor}
\usepackage{graphicx}
\usepackage{tabularx}
\usepackage{booktabs}
\usepackage{colortbl}
\usepackage{array}
\usepackage{enumitem}
\usepackage{fancyhdr}
\usepackage{titlesec}
\usepackage{lastpage}
\usepackage{setspace}
\usepackage{listings}
\usepackage{tcolorbox}

% Tipografía institucional Icesi (Sans-Serif Helvetica)
\usepackage{helvet}
\renewcommand{\familydefault}{\sfdefault}

% Interlineado aumentado un 30% (setstretch 1.38) para mayor holgura y lectura óptima
\setstretch{1.38}

% Altura de celdas en tablas
\renewcommand{\arraystretch}{1.25}

% Paleta de Colores Oficial de los Programas Icesi
\definecolor{icesiBlue}{HTML}{5454E9}       % Azul vibrante institucional
\definecolor{icesiOrange}{HTML}{F0681A}     % Naranja acento
\definecolor{icesiGreen}{HTML}{4CB978}      % Verde acento
\definecolor{icesiPurple}{HTML}{865CF0}     % Púrpura acento
\definecolor{icesiDark}{HTML}{1A1A1A}       % Texto principal oscuro
\definecolor{icesiMuted}{HTML}{6B7280}      % Gris metadatos
\definecolor{icesiBorder}{HTML}{D1D5DB}     % Líneas y bordes
\definecolor{icesiRowLine}{HTML}{E5E7EB}    % Líneas de tabla
\definecolor{icesiLight}{HTML}{F9FAFB}      % Fondo suave
\definecolor{codeBackground}{HTML}{F3F4F6}  % Fondo código

% Configuración de Listings
\lstdefinestyle{bashstyle}{
    backgroundcolor=\color{codeBackground},
    commentstyle=\color{icesiMuted}\itshape,
    keywordstyle=\color{icesiBlue}\bfseries,
    basicstyle=\ttfamily\scriptsize,
    breaklines=true,
    frame=single,
    rulecolor=\color{icesiBorder},
    language=bash,
    showstringspaces=false
}

\lstdefinelanguage{JavaScript}{
  keywords={typeof, new, true, false, catch, function, return, null, switch, var, if, in, while, do, else, case, break, const, let, async, await, import, export, from, default, try, throw, class, extends, super, Boolean},
  keywordstyle=\color{icesiBlue}\bfseries,
  ndkeywords={document, window, localStorage, JSON, console, fetch, Promise, Response, Error, String, Date},
  ndkeywordstyle=\color{icesiPurple}\bfseries,
  sensitive=true,
  comment=[l]{//},
  morecomment=[s]{/*}{*/},
  commentstyle=\color{icesiMuted}\itshape,
  stringstyle=\color{icesiOrange},
  morestring=[b]',
  morestring=[b]"
}

\lstdefinestyle{jsstyle}{
    language=JavaScript,
    backgroundcolor=\color{codeBackground},
    commentstyle=\color{icesiMuted}\itshape,
    keywordstyle=\color{icesiBlue}\bfseries,
    stringstyle=\color{icesiOrange},
    basicstyle=\ttfamily\scriptsize,
    breaklines=true,
    frame=single,
    rulecolor=\color{icesiBorder},
    showstringspaces=false,
    tabsize=2
}

\lstdefinelanguage{HTML5}{
  morekeywords={DOCTYPE, html, head, meta, title, link, body, header, main, section, article, form, input, button, select, option, ul, li, span, div, h1, h2, script},
  sensitive=false,
  morecomment=[s]{<!--}{-->},
  morestring=[b]"
}

\lstdefinestyle{htmlstyle}{
    language=HTML5,
    backgroundcolor=\color{codeBackground},
    basicstyle=\ttfamily\scriptsize,
    keywordstyle=\color{icesiPurple}\bfseries,
    commentstyle=\color{icesiMuted}\itshape,
    stringstyle=\color{icesiGreen},
    breaklines=true,
    frame=single,
    rulecolor=\color{icesiBorder},
    showstringspaces=false,
    tabsize=2
}

\lstset{
  inputencoding=utf8,
  extendedchars=true,
  literate=%
    {á}{{\'a}}1
    {é}{{\'e}}1
    {í}{{\'i}}1
    {ó}{{\'o}}1
    {ú}{{\'u}}1
    {Á}{{\'A}}1
    {É}{{\'E}}1
    {Í}{{\'I}}1
    {Ó}{{\'O}}1
    {Ú}{{\'U}}1
    {ñ}{{\~n}}1
    {Ñ}{{\~N}}1
    {¿}{{?`}}1
    {¡}{{!`}}1
}

% =========================================================================
% ENCABEZADO Y PIE DE PÁGINA ESTILO PROGRAMA DE CURSO ICESI
% =========================================================================
\pagestyle{fancy}
\fancyhf{}
\fancyhead[L]{%
  \begin{minipage}[c]{2.8cm}
    \includegraphics[height=0.58cm]{logo-icesi.png}
  \end{minipage}%
  \hspace{0.15cm}%
  {\color{icesiBorder}\vrule width 0.8pt height 0.60cm}%
  \hspace{0.2cm}%
  \begin{minipage}[c]{10.5cm}
    \fontsize{7.5}{9.5}\selectfont\textbf{\color{icesiDark}UNIVERSIDAD ICESI} \\
    \fontsize{7}{9}\selectfont\textbf{\color{icesiDark}GUÍA DE TALLER PRÁCTICO FORMATIVO: COMPUTACIÓN EN INTERNET I} \\
    \fontsize{6.5}{8.5}\selectfont\color{icesiMuted}Período: 2026-2 $\bullet$ NRC: 12378 / 12377 $\bullet$ Modalidad: Equipos de 2 Estudiantes
  \end{minipage}%
}
\fancyhead[R]{\fontsize{7.5}{9.5}\selectfont\color{icesiMuted}Pág. \thepage\ de \pageref{LastPage}}

% Barra horizontal multicolor característica del programa de curso
\renewcommand{\headrule}{%
  \vspace{0.08cm}%
  \noindent%
  \textcolor{icesiBlue}{\rule{0.60\textwidth}{1.8pt}}%
  \textcolor{icesiOrange}{\rule{0.18\textwidth}{1.8pt}}%
  \textcolor{icesiGreen}{\rule{0.11\textwidth}{1.8pt}}%
  \textcolor{icesiPurple}{\rule{0.11\textwidth}{1.8pt}}%
}
\renewcommand{\footrulewidth}{0.4pt}
\fancyfoot[L]{\fontsize{7.5}{9.5}\selectfont\color{icesiMuted}Computación en Internet I (09810 - TIC) $\bullet$ Arquitectura de Componentes en Vanilla JS y Promesas}
\fancyfoot[R]{\fontsize{7.5}{9.5}\selectfont\color{icesiMuted}Universidad Icesi $\bullet$ DCSI}

% =========================================================================
% COMANDOS DE SECCIONES CON BADGE ESTILO PROGRAMA ICESI
% =========================================================================
\newcommand{\sectionBadge}[1]{%
  \refstepcounter{section}%
  \vspace{0.24cm}%
  \noindent%
  \colorbox{icesiBlue}{\color{white}\rule[-0.08cm]{0pt}{0.42cm}\hspace{0.16cm}\textbf{\normalsize \thesection}\hspace{0.16cm}}%
  \hspace{0.18cm}%
  {\fontsize{9.5}{12}\selectfont\textbf{\color{icesiDark}#1}}\\[-0.10cm]%
  \noindent{\color{icesiDark}\rule{\textwidth}{0.8pt}}\vspace{0.10cm}%
}

\titleformat{\subsection}
  {\color{icesiBlue}\normalfont\normalsize\bfseries}
  {\thesubsection}{0.4em}{}
\titlespacing*{\subsection}{0pt}{0.6ex plus 0.2ex minus .1ex}{0.2ex plus .1ex}

\newcolumntype{C}[1]{>{\centering\arraybackslash}p{#1}}
\newcolumntype{L}[1]{>{\raggedright\arraybackslash}p{#1}}
\newcolumntype{Y}{>{\centering\arraybackslash}X}

\begin{document}
\deactivatequoting

% =========================================================================
% TÍTULO Y PRESENTACIÓN INICIAL (PÁGINA 1)
% =========================================================================
\noindent
{\fontsize{8.5}{11}\selectfont\textbf{\color{icesiBlue}TALLER PRÁCTICO FORMATIVO: ARQUITECTURA WEB CLIENTE Y MANIPULACIÓN AVANZADA DEL DOM}}\\[0.06cm]
{\fontsize{13}{16}\selectfont\textbf{Arquitectura Modular de Componentes en JavaScript Puro y Consumo Asíncrono con Promesas}}\\[0.10cm]
{\color{icesiRowLine}\hrule height 0.6pt}\vspace{0.10cm}

\noindent
{\color{icesiBlue}\textbf{\rule{2.5pt}{8.5pt}}}~\textbf{\fontsize{8.5}{11}\selectfont FACULTAD DE INGENIERÍA, DISEÑO Y CIENCIAS APLICADAS} \hfill
{\color{icesiOrange}\textbf{\rule{2.5pt}{8.5pt}}}~\textbf{\fontsize{8.5}{11}\selectfont DEPARTAMENTO DE COMPUTACIÓN Y SISTEMAS INTELIGENTES}\\[0.06cm]
{\color{icesiRowLine}\hrule height 0.6pt}\vspace{0.10cm}

% =========================================================================
% TABLA DE INFORMACIÓN Y METADATOS (ESTILO PROGRAMA DE CURSO)
% =========================================================================
\noindent
{\color{icesiBlue}\textbf{\rule{2.5pt}{8.5pt}}}~\textbf{\fontsize{8.5}{11}\selectfont INFORMACIÓN DE LA ACTIVIDAD} \hfill {\fontsize{7.5}{9}\selectfont\color{icesiMuted}\textbf{IDENTIFICACIÓN}}\\[-0.22cm]
{\color{icesiBorder}\hrule height 0.5pt}\vspace{0.05cm}

\noindent
\small
\begin{tabularx}{\textwidth}{>{\bfseries\color{icesiDark}\fontsize{7.5}{9.5}\selectfont}L{2.6cm}X!{\color{icesiBlue}\vrule width 0.8pt}>{\bfseries\color{icesiDark}\fontsize{7.5}{9.5}\selectfont}L{2.6cm}X}
CÓDIGO: & 09810 - TIC & PERÍODO: & 202620 (Semestre 2026-2) \\ \hline
\rowcolor{icesiLight}
NRC / GRUPO: & 12378 / 12377 (Grupos 001 y 003) & CARÁCTER: & Taller Práctico Formativo Guiado (Sin Calificación Sumativa) \\ \hline
MODALIDAD: & Equipos de 2 Estudiantes & DURACIÓN: & 2 Horas en Clase + Cierre Autónomo \\ \hline
\rowcolor{icesiLight}
POLÍTICA IAG: & Nivel 2: Asistencia de Consulta y Depuración & PONDERACIÓN: & Formativo (Alineación RAA 3 / Preparación Parcial 3) \\ \hline
ESTUDIANTE 1: & \rule{4.5cm}{0.4pt} & CÓDIGO: & \rule{2.2cm}{0.4pt} \\ \hline
\rowcolor{icesiLight}
ESTUDIANTE 2: & \rule{4.5cm}{0.4pt} & CÓDIGO: & \rule{2.2cm}{0.4pt} \\ \hline
\end{tabularx}

\vspace{0.10cm}

% =========================================================================
% SECCIÓN 1: CONTEXTO DEL PROBLEMA Y ARQUITECTURA DE COMPONENTES
% =========================================================================
\sectionBadge{Contexto: El Paradigma de Componentes en Vanilla JavaScript}

En las arquitecturas modernas de aplicaciones web de una sola página (\emph{Single Page Applications - SPAs}), el documento HTML inicial ya no contiene estructuras visuales fijas ni elementos estáticos como encabezados, formularios o tablas. En su lugar, el archivo \texttt{index.html} actúa exclusivamente como un \textbf{contenedor raíz vacío} (\texttt{<div id="root"></div>}) acompañado de la importación del script principal de JavaScript.

Toda la interfaz gráfica, la actualización dinámica y la manipulación del árbol DOM se construyen en tiempo de ejecución mediante código JavaScript puro. Para que este código sea mantenible, escalable y profesional, adoptamos el \textbf{Paradigma de Componentes Funcionales}:

\begin{itemize}[leftmargin=*,itemsep=1.5pt,topsep=2pt]
  \item \textbf{¿Qué es un Componente en Vanilla JS?:} Es una función modular reutilizable que recibe parámetros (conocidos como \textbf{props}) y retorna un nodo del DOM completamente configurado (con sus clases, atributos, eventos y contenido de texto seguro).
  \item \textbf{Principios de la Arquitectura Modular de Componentes:}
    \begin{itemize}
      \item \textbf{Props hacia abajo:} Los componentes padres pasan datos de configuración y estado hacia los componentes hijos mediante argumentos de función.
      \item \textbf{Event Callbacks hacia arriba:} Los componentes hijos notifican acciones del usuario (como presionar "Eliminar" o enviar el formulario) invocando funciones callback pasadas en sus props (ej. \texttt{onAddTask}, \texttt{onDeleteTask}).
      \item \textbf{Estado Centralizado (State):} El componente raíz (\texttt{App}) mantiene la única fuente de verdad de la aplicación (el array de tareas, el estado de carga y los mensajes de error). Cuando el estado cambia, se dispara una actualización dinámica de la interfaz.
    \end{itemize}
  \item \textbf{Consumo Asíncrono con Promesas (Service Layer):} El acceso a la red se encapsula en un módulo de servicio independiente (\texttt{taskService.js}). Cada operación HTTP (\texttt{GET}, \texttt{POST}, \texttt{DELETE}) retorna una \textbf{Promesa} que transita de forma determinística por sus tres estados:
    \begin{enumerate}[label=\alph*)]
      \item \textbf{Pending:} La petición está en vuelo; el componente \texttt{App} renderiza el componente \texttt{LoadingSpinner}.
      \item \textbf{Fulfilled:} La Promesa se resuelve con éxito; se actualiza la lista de tareas en el estado y se dibuja el componente \texttt{TaskList}.
      \item \textbf{Rejected:} La red o el servidor fallaron; se captura el motivo del rechazo en el bloque \texttt{catch} y se renderiza el componente \texttt{ErrorAlert}.
    \end{enumerate}
\end{itemize}

\begin{tcolorbox}[colback=icesiBlue!10,colframe=icesiBlue,arc=3pt,title=\textbf{Propósito Formativo del Taller}]
\fontsize{7.5}{9.5}\selectfont
Aprender a construir una aplicación web completa utilizando \textbf{Arquitectura de Componentes en Vanilla JS} y manipulación programática del DOM, partiendo de un HTML con una sola etiqueta (\texttt{<div id="root"></div>}). Desarrollaremos una suite modular de componentes (Header, TaskForm, TaskItem, TaskList, LoadingSpinner, ErrorAlert y App) y nos conectaremos a un servidor Express sencillo que implementa operaciones \texttt{GET}, \texttt{POST} y \texttt{DELETE} mediante \textbf{Promesas} con \texttt{fetch} y \texttt{async / await}.
\end{tcolorbox}

% =========================================================================
% SECCIÓN 2: LA API DE SOPORTE (CAJA NEGRA)
% =========================================================================
\sectionBadge{La API de Soporte: Servidor Express como Caja Negra}

Para abastecer nuestra aplicación frontend con datos reales, utilizaremos el servidor Express provisto por la cátedra. En este curso \textbf{no programamos el backend}; la API actúa como una caja negra en el puerto \texttt{3001} que maneja una lista de tareas en memoria.

\begin{figure}[h!]
\centering
\begin{lstlisting}[style=bashstyle]
  +---------------------------------------------------------------------------------------------------+
  |                                   FRONTEND: VANILLA JS COMPONENTS                                 |
  |                                       http://localhost:3000                                       |
  |                                                                                                   |
  |   index.html: <div id="root"></div>                                                               |
  |       |                                                                                           |
  |       v                                                                                           |
  |     App.js (Estado: tasks, loading, error)                                                        |
  |       |                                                                                           |
  |       +---> Header.js          +---> LoadingSpinner.js                                            |
  |       +---> TaskForm.js        +---> ErrorAlert.js                                                |
  |       +---> TaskList.js ----> TaskItem.js (con boton Eliminar)                                    |
  |       |                                                                                           |
  |       +---> taskService.js (Consumo con Promesas: fetch + async/await)                            |
  +---------------------------------------------------------------------------------------------------+
                  |                                 |                                 |
        1. GET /api/tasks                 2. POST /api/tasks                3. DELETE /api/tasks/:id
        (Listar tareas)                   (Crear nueva tarea)               (Eliminar tarea por ID)
                  |                                 |                                 |
                  v                                 v                                 v
  +---------------------------------------------------------------------------------------------------+
  |                                   API EXPRESS EN PUERTO 3001                                      |
  |                             https://github.com/Compu-1/api-express-promesas-taskflow               |
  +---------------------------------------------------------------------------------------------------+
\end{lstlisting}
\end{figure}

\subsection{Código Completo del Servidor (server.js)}
A continuación se presenta el código del backend para su lectura y contexto:

\begin{lstlisting}[style=jsstyle]
import express from 'express';
import cors from 'cors';

const app = express();
const PORT = 3001;

app.use(cors());          // Habilita peticiones cruzadas desde el frontend en localhost:3000
app.use(express.json());  // Middleware para interpretar cuerpos JSON en POST

// Lista de tareas inicial en memoria
let tareas = [
  { id: 1, title: "Comprender la estructura basica de HTML y CSS", completed: true },
  { id: 2, title: "Aprender a seleccionar y modificar elementos del DOM", completed: true },
  { id: 3, title: "Entender que es una Promesa y sus tres estados en JS", completed: false },
  { id: 4, title: "Consumir una API externa usando fetch y async/await", completed: false },
  { id: 5, title: "Renderizar dinamicamente los datos en la interfaz web", completed: false }
];

// 1. GET /api/tasks: Obtener todas las tareas
app.get('/api/tasks', (req, res) => res.json(tareas));

// 2. POST /api/tasks: Registrar nueva tarea
app.post('/api/tasks', (req, res) => {
  const { title } = req.body;
  if (!title || !title.trim()) return res.status(400).json({ error: "Titulo requerido" });
  const nuevaTarea = { id: Date.now(), title: title.trim(), completed: false };
  tareas.push(nuevaTarea);
  res.status(201).json(nuevaTarea);
});

// 3. DELETE /api/tasks/:id: Eliminar tarea por ID
app.delete('/api/tasks/:id', (req, res) => {
  const id = req.params.id;
  tareas = tareas.filter((t) => String(t.id) !== String(id));
  res.json({ message: "Tarea eliminada exitosamente", id });
});

app.listen(PORT, () => console.log(`Servidor activo en http://localhost:${PORT}`));
\end{lstlisting}

% =========================================================================
% SECCIÓN 3: PROCEDIMIENTO EXPERIMENTAL PASO A PASO
% =========================================================================
\sectionBadge{Procedimiento Experimental Paso a Paso: Construyendo el Front}

A continuación desarrollaremos la aplicación web estructurada modularmente en componentes funcionales.

\subsection{Fase 1: Puesta en Marcha de la API de Soporte (5 Minutos)}
En una terminal, clone e inicie el servidor de la cátedra:

\begin{lstlisting}[style=bashstyle]
git clone https://github.com/Compu-1/api-express-promesas-taskflow.git
cd api-express-promesas-taskflow
npm install && npm start
\end{lstlisting}

Verifique en su navegador que \texttt{http://localhost:3001/api/tasks} responde con el JSON de tareas. Mantenga esta terminal abierta.

\subsection{Fase 2: Estructura del Proyecto y el HTML Mínimo (10 Minutos)}
Abra una \textbf{segunda terminal} en una ubicación independiente para crear el proyecto cliente:

\begin{lstlisting}[style=bashstyle]
mkdir taskflow-components && cd taskflow-components
mkdir -p css js/services js/components
touch index.html css/styles.css js/main.js js/services/taskService.js
touch js/components/Header.js js/components/TaskForm.js js/components/TaskItem.js
touch js/components/TaskList.js js/components/LoadingSpinner.js js/components/ErrorAlert.js js/components/App.js
\end{lstlisting}

Edite \texttt{index.html}. Observe que el cuerpo del documento \textbf{solo tiene la etiqueta raíz \texttt{<div id="root"></div>}} y el enlace al script:

\begin{lstlisting}[style=htmlstyle]
<!DOCTYPE html>
<html lang="es">
<head>
  <meta charset="UTF-8">
  <meta name="viewport" content="width=device-width, initial-scale=1.0">
  <title>TaskFlow - Arquitectura de Componentes en Vanilla JS</title>
  <link rel="stylesheet" href="css/styles.css">
</head>
<body>
  <!-- Contenedor Raiz Unico: Toda la UI se crea dinamicamente via JS -->
  <div id="root"></div>

  <!-- Punto de entrada de la aplicacion -->
  <script type="module" src="js/main.js"></script>
</body>
</html>
\end{lstlisting}

Agregue los estilos visuales en \texttt{css/styles.css}:

\begin{lstlisting}[style=bashstyle]
/* css/styles.css: Paleta moderna institucional y estados de componentes */
:root { --primario: #5454E9; --exito: #4CB978; --error: #EF4444; --fondo: #F3F4F6; }
body { font-family: -apple-system, BlinkMacSystemFont, "Segoe UI", Roboto, sans-serif; background: var(--fondo); display: flex; justify-content: center; padding: 2.5rem 1rem; margin: 0; }
#root { width: 100%; max-width: 560px; }
.app-container { background: white; padding: 2.2rem; border-radius: 12px; box-shadow: 0 4px 18px rgba(0,0,0,0.08); }
.app-header h1 { margin: 0 0 0.3rem 0; color: #1F2937; font-size: 1.55rem; }
.app-subtitle { margin: 0; color: #6B7280; font-size: 0.88rem; }
.task-form { display: flex; gap: 0.5rem; margin: 1.5rem 0 1rem 0; }
.task-form input { flex: 1; padding: 0.7rem 0.9rem; border: 1.5px solid #D1D5DB; border-radius: 8px; font-size: 0.9rem; }
.task-form input:focus { outline: none; border-color: var(--primario); }
.btn-add { background: var(--exito); color: white; border: none; padding: 0.7rem 1.2rem; border-radius: 8px; font-weight: 600; cursor: pointer; white-space: nowrap; }
.btn-add:hover { background: #3EA668; }
.btn-refresh { background: var(--primario); color: white; border: none; padding: 0.65rem 1.2rem; border-radius: 8px; font-weight: 600; font-size: 0.9rem; cursor: pointer; width: 100%; margin-bottom: 1rem; }
.btn-refresh:hover { background: #3E3EC8; }
.btn-add:disabled, .btn-refresh:disabled { background: #9CA3AF; cursor: not-allowed; }
.loading-container { display: flex; align-items: center; justify-content: center; gap: 0.6rem; padding: 0.8rem 0; color: #4B5563; font-size: 0.88rem; }
.spinner { width: 18px; height: 18px; border: 2.5px solid #E5E7EB; border-top-color: var(--primario); border-radius: 50%; animation: girar 0.8s linear infinite; }
@keyframes girar { to { transform: rotate(360deg); } }
.error-alert { background: #FEE2E2; color: #991B1B; padding: 0.85rem; border-radius: 8px; font-size: 0.88rem; border: 1px solid #FCA5A5; margin-bottom: 1rem; display: flex; justify-content: space-between; align-items: center; gap: 0.5rem; }
.btn-retry { background: #DC2626; color: white; border: none; padding: 0.3rem 0.7rem; border-radius: 6px; font-size: 0.8rem; cursor: pointer; }
.task-list { list-style: none; padding: 0; margin: 0; display: flex; flex-direction: column; gap: 0.5rem; }
.task-item { display: flex; align-items: center; justify-content: space-between; padding: 0.75rem 0.9rem; background: #F9FAFB; border: 1px solid #E5E7EB; border-radius: 8px; gap: 0.8rem; }
.task-title { flex: 1; color: #1F2937; }
.task-item.completed .task-title { text-decoration: line-through; color: #9CA3AF; }
.task-actions { display: flex; align-items: center; gap: 0.5rem; }
.badge { font-size: 0.72rem; padding: 0.2rem 0.45rem; border-radius: 12px; font-weight: 600; white-space: nowrap; }
.badge-completed { background: #D1FAE5; color: #065F46; }
.badge-pending { background: #FEF3C7; color: #92400E; }
.btn-delete { background: transparent; border: 1px solid #FCA5A5; color: var(--error); border-radius: 6px; padding: 0.25rem 0.55rem; font-size: 0.76rem; cursor: pointer; font-weight: 600; }
.btn-delete:hover { background: #FEE2E2; }
.empty-message { text-align: center; color: #9CA3AF; padding: 1.5rem 0; }
\end{lstlisting}

\subsection{Fase 3: Módulo de Servicio con Promesas (js/services/taskService.js) (15 Minutos)}
Encapsulamos las llamadas \texttt{fetch} que retornan Promesas en un módulo reutilizable:

\begin{lstlisting}[style=jsstyle]
// js/services/taskService.js: Capa de comunicacion asincrona con la API
const API_URL = 'http://localhost:3001/api/tasks';

/**
 * 1. GET: Retorna una Promesa que resuelve en el array de tareas
 */
export async function getTasks() {
  const response = await fetch(API_URL);
  if (!response.ok) {
    throw new Error(`Fallo al consultar tareas. Codigo HTTP: ${response.status}`);
  }
  return await response.json(); // Segunda Promesa: deserializa JSON
}

/**
 * 2. POST: Retorna una Promesa que resuelve en el nuevo objeto tarea
 */
export async function createTask(title) {
  const response = await fetch(API_URL, {
    method: 'POST',
    headers: { 'Content-Type': 'application/json' },
    body: JSON.stringify({ title: title })
  });
  if (!response.ok) {
    throw new Error(`Fallo al crear la tarea. Codigo HTTP: ${response.status}`);
  }
  return await response.json();
}

/**
 * 3. DELETE: Retorna una Promesa que confirma la eliminacion
 */
export async function deleteTask(id) {
  const response = await fetch(`${API_URL}/${id}`, {
    method: 'DELETE'
  });
  if (!response.ok) {
    throw new Error(`Fallo al eliminar la tarea. Codigo HTTP: ${response.status}`);
  }
  return await response.json();
}
\end{lstlisting}

\subsection{Fase 4: Componentes Funcionales de Presentación (30 Minutos)}
Construiremos cada componente como una función pura que recibe \texttt{props} y retorna un nodo del DOM manipulado de forma segura con \texttt{createElement} y \texttt{textContent}.

\subsubsection{Componente 1: Header (js/components/Header.js)}
\begin{lstlisting}[style=jsstyle]
export function Header({ title, subtitle }) {
  const header = document.createElement('header');
  header.className = 'app-header';

  const h1 = document.createElement('h1');
  h1.textContent = title;

  const p = document.createElement('p');
  p.className = 'app-subtitle';
  p.textContent = subtitle;

  header.append(h1, p);
  return header;
}
\end{lstlisting}

\subsubsection{Componente 2: TaskForm (js/components/TaskForm.js)}
Emite un evento callback \texttt{onAddTask} al enviar el formulario:

\begin{lstlisting}[style=jsstyle]
export function TaskForm({ onAddTask, disabled }) {
  const form = document.createElement('form');
  form.className = 'task-form';

  const input = document.createElement('input');
  input.type = 'text';
  input.placeholder = 'Escribe una nueva tarea...';
  input.required = true;
  input.disabled = Boolean(disabled);

  const button = document.createElement('button');
  button.type = 'submit';
  button.className = 'btn-add';
  button.textContent = 'Agregar Tarea';
  button.disabled = Boolean(disabled);

  form.append(input, button);

  form.addEventListener('submit', (e) => {
    e.preventDefault();
    const value = input.value.trim();
    if (!value) return;
    onAddTask(value); // Callback al componente padre (App)
    input.value = '';
  });

  return form;
}
\end{lstlisting}

\subsubsection{Componente 3: TaskItem (js/components/TaskItem.js)}
Recibe una tarea y el callback \texttt{onDeleteTask}:

\begin{lstlisting}[style=jsstyle]
export function TaskItem({ task, onDeleteTask }) {
  const li = document.createElement('li');
  li.className = `task-item ${task.completed ? 'completed' : ''}`;

  const span = document.createElement('span');
  span.className = 'task-title';
  span.textContent = task.title;

  const actions = document.createElement('div');
  actions.className = 'task-actions';

  const badge = document.createElement('span');
  badge.className = `badge ${task.completed ? 'badge-completed' : 'badge-pending'}`;
  badge.textContent = task.completed ? 'Completada' : 'Pendiente';

  const btnDelete = document.createElement('button');
  btnDelete.className = 'btn-delete';
  btnDelete.textContent = 'Eliminar';
  btnDelete.onclick = () => onDeleteTask(task.id); // Callback al padre

  actions.append(badge, btnDelete);
  li.append(span, actions);
  return li;
}
\end{lstlisting}

\subsubsection{Componente 4: TaskList (js/components/TaskList.js)}
Recorre el array de tareas mapeando cada elemento a un \texttt{TaskItem}:

\begin{lstlisting}[style=jsstyle]
import { TaskItem } from './TaskItem.js';

export function TaskList({ tasks, onDeleteTask }) {
  const ul = document.createElement('ul');
  ul.className = 'task-list';

  if (!tasks || tasks.length === 0) {
    const empty = document.createElement('li');
    empty.className = 'empty-message';
    empty.textContent = 'No hay tareas registradas en la API.';
    ul.appendChild(empty);
    return ul;
  }

  tasks.forEach((task) => {
    ul.appendChild(TaskItem({ task, onDeleteTask }));
  });

  return ul;
}
\end{lstlisting}

\subsubsection{Componente 5 y 6: LoadingSpinner y ErrorAlert}
\begin{lstlisting}[style=jsstyle]
// js/components/LoadingSpinner.js: Muestra el estado PENDING de la Promesa
export function LoadingSpinner({ message = 'Consultando API remota...' }) {
  const container = document.createElement('div');
  container.className = 'loading-container';

  const spinner = document.createElement('div');
  spinner.className = 'spinner';

  const span = document.createElement('span');
  span.textContent = message;

  container.append(spinner, span);
  return container;
}
\end{lstlisting}

\begin{lstlisting}[style=jsstyle]
// js/components/ErrorAlert.js: Muestra el estado REJECTED de la Promesa
export function ErrorAlert({ message, onRetry }) {
  const alert = document.createElement('div');
  alert.className = 'error-alert';

  const text = document.createElement('span');
  text.textContent = message;
  alert.appendChild(text);

  if (onRetry) {
    const btnRetry = document.createElement('button');
    btnRetry.className = 'btn-retry';
    btnRetry.textContent = 'Reintentar';
    btnRetry.onclick = onRetry;
    alert.appendChild(btnRetry);
  }

  return alert;
}
\end{lstlisting}

\subsection{Fase 5: Componente Raíz y Orquestador de Estado (js/components/App.js) (25 Minutos)}
El componente \texttt{App} administra el estado centralizado, reacciona a los eventos de los hijos y consume el servicio asíncrono controlando los estados de las Promesas:

\begin{lstlisting}[style=jsstyle]
// js/components/App.js: Componente Raiz con Estado y Orquestacion
import { getTasks, createTask, deleteTask } from '../services/taskService.js';
import { Header } from './Header.js';
import { TaskForm } from './TaskForm.js';
import { TaskList } from './TaskList.js';
import { LoadingSpinner } from './LoadingSpinner.js';
import { ErrorAlert } from './ErrorAlert.js';

export function App() {
  const container = document.createElement('main');
  container.className = 'app-container';

  // 1. Estado centralizado de la aplicacion
  const state = {
    tasks: [],
    loading: false,
    loadingMessage: '',
    error: null
  };

  // 2. Funcion de renderizado dinamico (construye el DOM segun el estado actual)
  function render() {
    container.innerHTML = ''; // Limpia el arbol anterior para redibujar

    // Encabezado
    container.appendChild(Header({
      title: 'TaskFlow Components',
      subtitle: 'Arquitectura de Componentes en Vanilla JS y Promesas (GET, POST, DELETE)'
    }));

    // Formulario de Creacion (POST)
    container.appendChild(TaskForm({
      disabled: state.loading,
      onAddTask: handleAddTask
    }));

    // Indicador de Carga (Estado PENDING)
    if (state.loading) {
      container.appendChild(LoadingSpinner({ message: state.loadingMessage }));
    }

    // Alerta de Error (Estado REJECTED)
    if (state.error) {
      container.appendChild(ErrorAlert({
        message: state.error,
        onRetry: handleLoadTasks
      }));
    }

    // Boton de refresco manual
    const refreshBtn = document.createElement('button');
    refreshBtn.className = 'btn-refresh';
    refreshBtn.textContent = 'Refrescar Lista desde API';
    refreshBtn.disabled = state.loading;
    refreshBtn.onclick = handleLoadTasks;
    container.appendChild(refreshBtn);

    // Lista de Tareas (Estado FULFILLED)
    container.appendChild(TaskList({
      tasks: state.tasks,
      onDeleteTask: handleDeleteTask
    }));
  }

  // 3. Manejadores Asincronos de Eventos (Consumo de Promesas)
  async function handleLoadTasks() {
    state.loading = true;
    state.loadingMessage = 'Consultando API remota... (Estado: PENDING)';
    state.error = null;
    render();

    try {
      state.tasks = await getTasks(); // Promesa GET resuelta (FULFILLED)
    } catch (err) {
      state.error = `Error al consultar la API: ${err.message}. Verifica el servidor en el puerto 3001.`;
    } finally {
      state.loading = false;
      render();
    }
  }

  async function handleAddTask(title) {
    state.loading = true;
    state.loadingMessage = `Guardando "${title}" en la API... (Estado: PENDING)`;
    state.error = null;
    render();

    try {
      const nueva = await createTask(title); // Promesa POST resuelta (FULFILLED)
      state.tasks.push(nueva);
    } catch (err) {
      state.error = `Error al crear tarea: ${err.message}`;
    } finally {
      state.loading = false;
      render();
    }
  }

  async function handleDeleteTask(id) {
    state.loading = true;
    state.loadingMessage = `Eliminando tarea ID ${id}... (Estado: PENDING)`;
    state.error = null;
    render();

    try {
      await deleteTask(id); // Promesa DELETE resuelta (FULFILLED)
      state.tasks = state.tasks.filter((t) => String(t.id) !== String(id));
    } catch (err) {
      state.error = `Error al eliminar tarea: ${err.message}`;
    } finally {
      state.loading = false;
      render();
    }
  }

  // Carga inicial automatica de datos
  handleLoadTasks();

  return container;
}
\end{lstlisting}

\subsection{Fase 6: Montaje en el Punto de Entrada (js/main.js) (5 Minutos)}
En \texttt{js/main.js}, importamos \texttt{App} y la insertamos en el contenedor raíz \texttt{\#root}:

\begin{lstlisting}[style=jsstyle]
// js/main.js: Punto de entrada de la aplicacion
import { App } from './components/App.js';

const root = document.querySelector('#root');
root.appendChild(App());
\end{lstlisting}

\subsection{Fase 7: Ejecución y Experimentación en Vivo (15 Minutos)}
Inicie el servidor estático del frontend en \texttt{http://localhost:3000}:

\begin{lstlisting}[style=bashstyle]
npx live-server --port=3000
\end{lstlisting}

Realice los siguientes **4 experimentos guiados en vivo**:

\begin{enumerate}[leftmargin=*,itemsep=3pt,topsep=2pt]
  \item \textbf{Experimento 1: Inspección del DOM Vacío vs. Renderizado Dinámico}
    \begin{itemize}
      \item En su navegador, abra DevTools (\texttt{F12}) y seleccione la pestaña \textbf{Elements (Elementos)}.
      \item Abra el archivo fuente original \texttt{index.html} y compruebe que solo contiene \texttt{<div id="root"></div>}.
      \item Observe en el árbol del inspector cómo JavaScript construyó toda la jerarquía de elementos (\texttt{header}, \texttt{form}, \texttt{button}, \texttt{ul}, \texttt{li}) dentro de \texttt{\#root}.
    \end{itemize}
  \item \textbf{Experimento 2: Ciclo de Vida de Promesa en Creación (POST)}
    \begin{itemize}
      \item Escriba una tarea en el input del componente \texttt{TaskForm} y haga clic en "Agregar Tarea".
      \item Observe cómo el componente \texttt{App} cambia inmediatamente su estado a \texttt{loading: true}, renderizando el \texttt{LoadingSpinner} (\textbf{Estado Pending}).
      \item Al resolver la Promesa, la API entrega el ID asignado, el componente añade la tarea a \texttt{state.tasks} y re-renderiza la lista (\textbf{Estado Fulfilled}).
    \end{itemize}
  \item \textbf{Experimento 3: Flujo de Event Callbacks en Eliminación (DELETE)}
    \begin{itemize}
      \item Haga clic en el botón "Eliminar" de un elemento de la lista.
      \item Observe cómo el componente hijo \texttt{TaskItem} notifica el evento al padre \texttt{App} mediante \texttt{onDeleteTask(id)}.
      \item Al completarse la Promesa \texttt{DELETE}, el item desaparece de la vista sin recargar la página.
    \end{itemize}
  \item \textbf{Experimento 4: Estado REJECTED y Resiliencia con ErrorAlert}
    \begin{itemize}
      \item En la terminal de Express, detenga el servidor con \texttt{Ctrl + C}.
      \item En el navegador, intente agregar una tarea o presione el botón de refresco.
      \item Compruebe que la Promesa de red falla, la ejecución salta al bloque \texttt{catch} y el componente \texttt{ErrorAlert} se dibuja con el botón de "Reintentar".
      \item Reinicie el servidor con \texttt{npm start}, presione "Reintentar" y verifique la recuperación transparente de la aplicación.
    \end{itemize}
\end{enumerate}

% =========================================================================
% SECCIÓN 4: CUESTIONARIO DE ANÁLISIS CONCEPTUAL
% =========================================================================
\sectionBadge{Cuestionario de Análisis Conceptual y Discusión Técnica}

Responda en pareja las siguientes preguntas conceptuales:

\begin{enumerate}[leftmargin=*,itemsep=4pt,topsep=2pt]
  \item \textbf{Arquitectura de Componentes frente a HTML tradicional:} ¿Cuáles son las ventajas de tener un archivo \texttt{index.html} que solo contenga \texttt{<div id="root"></div>} y generar la interfaz mediante componentes funcionales en JavaScript? ¿Cómo facilita esto la modularidad, depuración y mantenimiento del código frente a un único script monolítico?
  \item \textbf{Flujo Unidireccional de Datos (Props y Callbacks):} Explique cómo viajan los datos y las acciones entre el componente \texttt{App}, \texttt{TaskList} y \texttt{TaskItem}:
    \begin{itemize}
      \item ¿Cómo recibe \texttt{TaskItem} la información de la tarea que debe pintar?
      \item ¿Cómo le avisa \texttt{TaskItem} al componente \texttt{App} que el usuario presionó el botón de eliminar?
    \end{itemize}
  \item \textbf{Separación de Responsabilidades (Service Layer vs. Componentes):} ¿Por qué es una buena práctica de ingeniería extraer las funciones \texttt{fetch} al archivo \texttt{taskService.js} en lugar de escribir \texttt{fetch()} directamente dentro de cada componente visual?
  \item \textbf{Manejo de Estados Asíncronos en la UI:} Describa qué componentes visuales se montan o desmontan del DOM según los tres estados de la Promesa:
    \begin{itemize}
      \item En estado \textbf{Pending}:
      \item En estado \textbf{Fulfilled}:
      \item En estado \textbf{Rejected}:
    \end{itemize}
\end{enumerate}

% =========================================================================
% SECCIÓN 5: RÚBRICA DE EVALUACIÓN MODULARIZADA
% =========================================================================
\sectionBadge{Rúbrica de Autoevaluación Formativa Guiada (Escala 0.0 a 5.0)}

\begin{table}[h!]
\centering
\fontsize{7}{9}\selectfont
\begin{tabularx}{\textwidth}{L{2.5cm}X X X C{1.1cm}}
\toprule
\textbf{Criterio Evaluado} & \textbf{Excelente (4.5 -- 5.0)} & \textbf{Aceptable (3.0 -- 4.4)} & \textbf{Insuficiente (0.0 -- 2.9)} & \textbf{Puntaje} \\
\midrule
\textbf{1. Arquitectura de Componentes y DOM} &
Estructura la aplicación en componentes funcionales modulares a partir de un único \texttt{<div id="root"></div>}, manipulando el DOM con \texttt{createElement} y \texttt{textContent}. &
Implementa componentes pero recurre a estructuras HTML estáticas fuera de \texttt{\#root} o abusa de \texttt{innerHTML}. &
No modulariza en componentes o mantiene toda la vista codificada rígidamente en el archivo HTML. &
\textbf{1.5 pts} \\
\rowcolor{icesiLight}
\textbf{2. Consumo de Promesas (Service Layer)} &
Implementa el servicio \texttt{taskService.js} con \texttt{fetch} y \texttt{async/await} para \texttt{GET}, \texttt{POST} y \texttt{DELETE}, gestionando \texttt{try/catch/finally}. &
Consume las Promesas pero escribe las llamadas directamente en los componentes sin capa de servicio. &
Errores de sintaxis en el consumo de red o incapacidad de procesar las respuestas de la API. &
\textbf{1.5 pts} \\
\textbf{3. Manejo de Estado y Eventos (Props / Callbacks)} &
El componente \texttt{App} centraliza el estado y orquesta el renderizado dinámico, pasando props y manejando callbacks correctamente. &
Maneja el estado pero la actualización visual del DOM presenta inconsistencias al crear o eliminar elementos. &
La interfaz no responde a eventos o colapsa ante cambios en los datos. &
\textbf{1.0 pt} \\
\rowcolor{icesiLight}
\textbf{4. Discusión Conceptual y Resiliencia} &
Explica con precisión el flujo unidireccional de datos y demuestra la captura de errores en \texttt{ErrorAlert} al apagar la API. &
Comprende la estructura de componentes pero tiene dudas sobre la relación entre props y callbacks. &
No responde el cuestionario o desconoce el funcionamiento asíncrono de las Promesas. &
\textbf{1.0 pt} \\
\bottomrule
\multicolumn{4}{r}{\textbf{PUNTAJE FORMATIVO TOTAL:}} & \textbf{5.0 pts} \\
\bottomrule
\end{tabularx}
\end{table}

% =========================================================================
% SECCIÓN 6: INSTRUCCIONES DE VERIFICACIÓN
% =========================================================================
\sectionBadge{Instrucciones de Verificación y Cierre en Aula}

Al finalizar la sesión práctica, presente ante el profesor o monitor:
\begin{enumerate}[leftmargin=*,itemsep=2pt,topsep=2pt]
  \item Inspección en DevTools mostrando que el archivo \texttt{index.html} solo tiene \texttt{<div id="root"></div>} y que toda la UI fue construida mediante componentes JavaScript.
  \item Demostración de creación de tarea (\texttt{POST}) y eliminación (\texttt{DELETE}) con actualización dinámica de la interfaz.
  \item Demostración de resiliencia: apagar la API y mostrar la aparición del componente \texttt{ErrorAlert} con la opción de reintentar.
\end{enumerate}

\end{document}
